% ============================================================
% Project: THE PHILOSOPHY OF VIABILITY
% Book: Book I — Ontology of Viability
% Book ID: VS.B01
% Part: Part I — Identity and Manifestation
% Chapter: Chapter 07 — Formal and Logical Reconstruction
% Section: 04 — Definitions and Conceptual Dependencies
% Asset: TL01 — Definition and Dependency Audit
% File: VS.B01.P01.C07.S04.SS04.P001.TL01.definition_and_dependency_audit.tex
% Purpose: To classify definitional operations and dependency
% relations for systematic conceptual and philosophical review.
% Placement: Section 04, subsection
% "A Method for Constructing and Auditing Definitions".
% ============================================================

\begin{table}[htbp]
\centering
\small
\setlength{\tabcolsep}{5pt}
\renewcommand{\arraystretch}{1.35}
\begin{tabular}{|p{2.7cm}|p{4.0cm}|p{4.1cm}|p{3.0cm}|}
\hline
\textbf{Category} &
\textbf{What it establishes} &
\textbf{What it does not establish by itself} &
\textbf{Audit question} \\
\hline

Stipulative definition &
A rule assigning a specified use to an expression within a
framework. &
The independent reality or ontological truth of the stipulated
concept. &
Is the scope of the stipulation explicit? \\
\hline

Explication &
A clarified and more precise articulation of a concept. &
That the proposed clarification is uniquely correct or
ontologically exhaustive. &
Does it preserve the relevant conceptual problem? \\
\hline

Explicit definition &
A stated condition or equivalence governing an expression. &
The ontological adequacy of the expression or its interpretation. &
Are the defining terms independently intelligible? \\
\hline

Implicit characterization &
A conceptual or formal role specified through a system of
relations or constraints. &
A unique interpretation unless uniqueness is separately shown. &
Can alternative interpretations satisfy the same constraints? \\
\hline

Conceptual dependence &
A relation of articulation or intelligibility between concepts. &
Ontological dependence or logical entailment. &
What conceptual presupposition is involved? \\
\hline

Logical dependence &
A relation of consequence between propositions under specified
premises and rules. &
The truth or justification of the premises. &
Which inference rules and premises are required? \\
\hline

Ontological dependence &
A substantive relation of grounding, constitution, realization,
or other specified ontological dependence. &
A merely conceptual or expository priority. &
What kind of ontological dependence is asserted? \\
\hline

Epistemic or justificatory dependence &
A relation in which the warrant for one claim relies on another. &
An ontological relation between the matters described. &
What evidence or reason supports the dependent claim? \\
\hline

Conceptual primitive &
A concept left undefined within a specified conceptual
reconstruction. &
Absolute simplicity or ontological fundamentality. &
Is its role sufficiently clarified for the present purpose? \\
\hline

Formal primitive &
An expression left undefined within a formal language. &
The philosophical primitiveness of the represented concept. &
What interpretation is assigned to the expression? \\
\hline

\end{tabular}
\caption{Classification of definitions and dependency relations
for the audit of conceptual and formal reconstruction.}
\label{tab:definition-and-dependency-audit}
\end{table}

% ============================================================
% END OF ASSET TL01 — CHAPTER 07, SECTION 04
% ============================================================